
%%%%%%%%%%%%%%%%%%%%%%%%%%%%%%%%%%%%%%%%
%           main body
%%%%%%%%%%%%%%%%%%%%%%%%%%%%%%%%%%%%%%%%

%-----------------------------------
%	TITLE PAGE
%-----------------------------------

\title[第七章\\
定积分]{\texorpdfstring{\LARGE \bfseries 第七章\quad 定积分}{第七章 定积分}} % The short title appears at the bottom of every slide, the full title is only on the title page
\author[\smaller \S 7.1\\
定积分概念和可积条件] % Your institution as it will appear on the bottom of every slide, maybe shorthand to save space
{\Large \bfseries \S 7.1 定积分的概念和可积条件}
% \author{Singular Value Decomposition} % Your name
% \date{\today} % Date, can be changed to a custom date
\date{}

%------------------------------------------------

\begin{document}

%------------------------------------------------

\begin{frame}

\titlepage % Print the title page as the first slide

\end{frame}

%------------------------------------------------

% \begin{frame}
% \frametitle{内容提要} % Table of contents slide, comment this block out to remove it
% \tableofcontents % Throughout your presentation, if you choose to use \section{} and \subsection{} commands, these will automatically be printed on this slide as an overview of your presentation
% \end{frame}

%------------------------------------------------
 
\section{定积分概念的导出背景}

%------------------------------------------------
 
% \begin{frame}
% \frametitle{历史背景与几何原型}
% \begin{block}{核心问题}
% 如何求曲边梯形的面积？
% \begin{itemize}
%     \item 直线图形面积 $\rightarrow$ 曲线图形面积
%     \item 有限分割：将曲边梯形分割为窄曲边梯形
%     \item 局部近似：用矩形面积近似代替窄曲边梯形
%     \item 极限过程：通过无限细分实现精确化
% \end{itemize}
% \end{block}
 
% \end{frame}

%------------------------------------------------
 
\begin{frame}

{\larger[2]
问题:

\vspace{1em}

\begin{itemize}[<+->]
\item 这个极限过程是什么?
\vspace{0.5em}
\item 是否总有极限?
\vspace{0.5em}
\item 有没有别的积分方式?
\end{itemize}
}
 
\end{frame}

%------------------------------------------------

\section{定积分的严格定义}

%------------------------------------------------

\begin{frame}
\frametitle{闭区间上的划分}

\begin{Def}[划分, partition]
设闭区间 $[a,b] \subset \mathbb{R},$ 称有限点集$P = \{x_0,x_1,\dots,x_n\}$ 为闭区间的一个\alert{划分}, 若满足:
\begin{equation*}
a = x_0 < x_1 < \cdots < x_{n-1} < x_n = b 
\end{equation*}
每个子区间 $[x_{k-1},x_k]$ 的长度记为 $\Delta x_k = x_k - x_{k-1},$ $k=1,2,\dots,n$
\end{Def}
 
\begin{Def}[带标志点的划分, Tagged Partition]
对任意划分 $P = \{x_0,x_1,\dots,x_n\},$ 若在每个子区间 $[x_{k-1},x_k]$ 中选取一点 $\xi_k \in [x_{k-1},x_k],$ 则称有序对
\[
(P,\xi) = (\{x_k\}_{k=0}^n,\ \{\xi_k\}_{k=1}^n)
\]
为\alert{带标志点的划分}，其中 $\xi = \{\xi_1,\dots,\xi_n\}$ 称为标志点集.
\end{Def}

\end{frame}
 
%------------------------------------------------

\begin{frame}
\frametitle{闭区间上的划分}

$$
\begin{array}{cccccccccccc}
P: ~~ & a=x_0 & < & x_1 & < & x_2 & < & \cdots & < & x_{n-1} & < & x_n = b \\
\xi: ~~ & & \xi_1 & & \xi_2 & & \xi_3 & \cdots & \xi_{n-1} & & \xi_n
\end{array}
$$

\pause

\begin{block}{划分相关的一些记号}
\begin{itemize}
\item 划分的\alert{模} (norm) 或\alert{细度} (mesh): $\lambda(P) = \max\limits_{1 \leq k \leq n} \Delta x_k$
\item 积分和 (黎曼和)：$S(f; P,\xi) = \sum_{k=1}^n f(\xi_k)\Delta x_k$
\end{itemize}
\end{block}

\pause

\begin{attnblock}
划分的精细程度由模 $\lambda(P)$ 控制, 当 $\lambda(P) \to 0$ 时所有子区间同步收缩, 这是黎曼可积的关键条件.
\end{attnblock}

\end{frame}
 
%------------------------------------------------

\begin{frame}
\frametitle{关于滤子基的极限}

\end{frame}
 
%------------------------------------------------
 
\begin{frame}
\frametitle{黎曼可积与黎曼积分}

\begin{Def}[函数黎曼可积及其黎曼积分]
设函数 $f$ 是定义在 $[a,b]$ 上的 (有界) 函数, 若关于 $\lambda(P) \to 0$ 的极限存在：
\begin{equation*}
\lim_{\lambda(P) \to 0} S(f; P,\xi) = \lim_{\lambda(P) \to 0} \sum_{k=1}^n f(\xi_k)\Delta x_k = I,
\end{equation*}
则称 $f$ 在 $[a,b]$ 上黎曼可积, 实数 $I$ 称为 $f$ 在 $[a, b]$ 上的黎曼积分 (定积分), 记作：
\begin{equation*}
I = \int_a^b f(x)dx
\end{equation*}

% \small
% \[
% \forall~\varepsilon > 0, ~ \exists~\delta>0, ~ \text{当分割细度} ~ \lambda(P) <\delta \Longrightarrow \left|\sum_{i=1}^n f(\xi_i)\Delta x_i - I\right| < \varepsilon
% \]
% \normalsize
% 则称 $f$ 在 $[a,b]$ 上黎曼可积, 
\end{Def}
 
% \begin{itemize}
% \item 分割的任意性：分点与取样点均任意选取
% \item 积分和的独立性：与分割方式和取样点选择无关
% \end{itemize}

\end{frame}

%------------------------------------------------

\section{达布和与可积条件}
 
%------------------------------------------------

\begin{frame}
\frametitle{达布上下和}
\begin{block}{关键概念}
\begin{itemize}
\item 分割 $P$ 的达布上和：$\overline{S}(P;f) = \sum_{i=1}^n M_i\Delta x_i$
\item 达布下和：$\underline{S}(P;f) = \sum_{i=1}^n m_i\Delta x_i$
\item 达布上、下积分：$\overline{\int}_a^b f = \inf_P U(P,f)$，$\underline{\int}_a^b f = \sup_P L(P,f)$
\end{itemize}
\end{block}
 
\begin{thm}[可积充要条件]
函数 $f$ 在 $[a,b]$ 上可积的充要条件是：
$$
\forall \varepsilon > 0,~\exists~\text{分割}~P,\ 使得~\overline{S}(P,f) - \underline{S}(P,f) < \varepsilon
$$
\end{thm}
\end{frame}

%-----------------------------------
%	PRESENTATION SLIDES
%-----------------------------------

\end{document}
